\documentclass{beamer}

\title[M.C.]{Markov Chains}
\subtitle{MATH 1554: \textit{Linear Algebra}}
\author{Ahmed Mohammed}
\date{PLUS Session, June 2023}
\institute{Georgia Institute of Technology}

\usetheme{Berlin}
\AtBeginSection[] {
	\begin{frame}
		\frametitle{Table Contents}
		\tableofcontents[currentsection]
	\end{frame}
}

\DeclareMathOperator{\vspan}{Span}

\begin{document}
\frame{\titlepage}

\section{Quick Intro Into Markov Chains}
\begin{frame}
	\frametitle{Opening Activity: Introduction}
	We wish to model a discrete dynamic system. For this, imagine a vector \(\vec{x}_0\) that describes the initial state of the system. It's a probability vector, so the entries
	add up to 1 and are non-negative.\pause
	
	Now, the system transitions between ``states" each step. \[\vec{x}_0 \to \vec{x}_1 \to \dots \vec{x}_k \to \dots\] \pause

	The sequence of probability vectors \(\{\vec{x}_n\}\) and the matrix \(P\) that controls the transition between states are collectively called the \textit{Markov Chain.}
\end{frame}

\begin{frame}
	\frametitle{Opening Activity: Introduction}
	The matrix \(P\) is called a \textit{stochastic matrix.} It's columns are also probability vectors. The characterstics of the system sometimes rely on \(P\). \pause
	Regular vs. non-regular. \pause

	For now, that doesn't matter. If \(P\) is a special type of matrix, we know that there exists one vector \(\vec{q}\) such that \[P\vec{q} = \vec{q}.\] \pause
	We stabalize into that state.
\end{frame}

\begin{frame}
	\frametitle{Main Activity: Solve a Problem}
	\begin{center}
		Refer to Worksheet 3!
	\end{center}

\end{frame}

\section{Review for Exam 2}
\begin{frame}
	\frametitle{Closing Activity: Review}
	\begin{center}
		Refer to worksheet 2!
	\end{center}
\end{frame}

\end{document}